
% BEGIN bib
 
\par\medskip\noindent\textbf{References}\par
  
\par\noindent  [1]  María Alpuente, Demis Ballis, Angel Cuenca-Ortega, Santiago Escobar, and José Meseguer.  ACUOS${}^{\mbox{2}}$: A High-Performance System for Modular ACU Generalization with Subtyping and Inheritance.  In Francesco Calimeri, Nicola Leone, and Marco Manna, editors, Logics in Artificial Intelligence - 16th European Conference, JELIA 2019, Rende, Italy, May 7-11, 2019, Proceedings, volume 11468 of Lecture Notes in Computer Science, pages 171–181. Springer, 2019.  doi:10.1007/978-3-030-19570-0\_11.  
 
\par\noindent  [2]  María Alpuente, Santiago Escobar, Javier Espert, and José Meseguer.  A modular order-sorted equational generalization algorithm.  Inf. Comput., 235:98–136, 2014.  doi:10.1016/J.IC.2014.01.006.  
 
\par\noindent  [3]  Nino Amiridze and Temur Kutsia.  Anti-unification and natural language processing.  EasyChair Preprint no. 203, EasyChair, 2018.  doi:10.29007/fkrh.  
 
\par\noindent  [4]  Mauricio Ayala-Rincón, David M. Cerna, Andres Felipe Gonzalez Barragan, and Temur Kutsia.  Equational anti-unification over absorption theories.  In Christoph Benzmüller, Marijn J. H. Heule, and Renate A. Schmidt, editors, Automated Reasoning - 12th International Joint Conference, IJCAR 2024, Nancy, France, July 3-6, 2024, Proceedings, Part II, volume 14740 of Lecture Notes in Computer Science, pages 317–337. Springer, 2024.  doi:10.1007/978-3-031-63501-4\_17.  
 
\par\noindent  [5]  Franz Baader.  Unification, weak unification, upper bound, lower bound, and generalization problems.  In Ronald V. Book, editor, Rewriting Techniques and Applications, 4th International Conference, RTA-91, Como, Italy, April 10-12, 1991, Proceedings, volume 488 of Lecture Notes in Computer Science, pages 86–97. Springer, 1991.  doi:10.1007/3-540-53904-2\_88.  
 
\par\noindent  [6]  Franz Baader and Tobias Nipkow.  Term rewriting and all that.  Cambridge University Press, 1998.  doi:10.1017/CBO9781139172752.  
 
\par\noindent  [7]  Franz Baader and Klaus U. Schulz.  Unification in the union of disjoint equational theories: Combining decision procedures.  J. Symb. Comput., 21(2):211–243, 1996.  doi:10.1006/JSCO.1996.0009.  
 
\par\noindent  [8]  Johannes Bader, Andrew Scott, Michael Pradel, and Satish Chandra.  Getafix: learning to fix bugs automatically.  Proc. ACM Program. Lang., 3(OOPSLA):159:1–159:27, 2019.  doi:10.1145/3360585.  
 
\par\noindent  [9]  Adam D. Barwell, Christopher Brown, and Kevin Hammond.  Finding parallel functional pearls: Automatic parallel recursion scheme detection in Haskell functions via anti-unification.  Future Gener. Comput. Syst., 79:669–686, 2018.  doi:10.1016/j.future.2017.07.024.  
 
\par\noindent  [10]  Hans-Jürgen Bürckert, Alexander Herold, and Manfred Schmidt-Schauß.  On equational theories, unification, and (un)decidability.  J. Symb. Comput., 8(1/2):3–49, 1989.  doi:10.1016/S0747-7171(89)80021-5.  
 
\par\noindent  [11]  Jochen Burghardt.  E-generalization using grammars.  Artif. Intell., 165(1):1–35, 2005.  doi:10.1016/J.ARTINT.2005.01.008.  
 
\par\noindent  [12]  David Cao, Rose Kunkel, Chandrakana Nandi, Max Willsey, Zachary Tatlock, and Nadia Polikarpova.  babble: Learning Better Abstractions with E-Graphs and Anti-unification.  Proceedings of the ACM on Programming Languages, 7(POPL):396–424, 2023.  doi:10.1145/3571207.  
 
\par\noindent  [13]  David M. Cerna.  Anti-unification and the theory of semirings.  Theor. Comput. Sci., 848:133–139, 2020.  doi:10.1016/J.TCS.2020.10.020.  
 
\par\noindent  [14]  David M. Cerna and Temur Kutsia.  Idempotent anti-unification.  ACM Trans. Comput. Log., 21(2):10:1–10:32, 2020.  doi:10.1145/3359060.  
 
\par\noindent  [15]  David M. Cerna and Temur Kutsia.  Unital anti-unification: Type and algorithms.  In Zena M. Ariola, editor, 5th International Conference on Formal Structures for Computation and Deduction, FSCD 2020, June 29-July 6, 2020, Paris, France (Virtual Conference), volume 167 of LIPIcs, pages 26:1–26:20. Schloss Dagstuhl – Leibniz-Zentrum für Informatik, 2020.  doi:10.4230/LIPICS.FSCD.2020.26.  
 
\par\noindent  [16]  David M. Cerna and Temur Kutsia.  Anti-unification and generalization: A survey.  In Proceedings of the Thirty-Second International Joint Conference on Artificial Intelligence, IJCAI 2023, 19th-25th August 2023, Macao, SAR, China, pages 6563–6573. ijcai.org, 2023.  doi:10.24963/IJCAI.2023/736.  
 
\par\noindent  [17]  Hubert Comon, Marianne Haberstrau, and Jean-Pierre Jouannaud.  Syntacticness, cycle-syntacticness, and shallow theories.  Inf. Comput., 111(1):154–191, 1994.  doi:10.1006/inco.1994.1043.  
 
\par\noindent  [18]  Andrew Cropper, Sebastijan Dumancic, Richard Evans, and Stephen H. Muggleton.  Inductive logic programming at 30.  Mach. Learn., 111(1):147–172, 2022.  doi:10.1007/S10994-021-06089-1.  
 
\par\noindent  [19]  Reudismam Rolim de Sousa, Gustavo Soares, Rohit Gheyi, Titus Barik, and Loris D’Antoni.  Learning quick fixes from code repositories.  In Cristiano D. Vasconcellos, Karina Girardi Roggia, Vanessa Collere, and Paulo Bousfield, editors, 35th Brazilian Symposium on Software Engineering, SBES 2021, Joinville, Santa Catarina, Brazil, 27 September 2021 - 1 October 2021, pages 74–83. ACM, 2021.  doi:10.1145/3474624.3474650.  
 
\par\noindent  [20]  Eric Domenjoud, Francis Klay, and Christophe Ringeissen.  Combination techniques for non-disjoint equational theories.  In Alan Bundy, editor, Automated Deduction - CADE-12, 12th International Conference on Automated Deduction, Nancy, France, June 26 - July 1, 1994, Proceedings, volume 814 of Lecture Notes in Computer Science, pages 267–281. Springer, 1994.  doi:10.1007/3-540-58156-1\_19.  
 
\par\noindent  [21]  Rui Dong, Zhicheng Huang, Ian Iong Lam, Yan Chen, and Xinyu Wang.  WebRobot: web robotic process automation using interactive programming-by-demonstration.  In Ranjit Jhala and Isil Dillig, editors, PLDI ’22: 43rd ACM SIGPLAN International Conference on Programming Language Design and Implementation, San Diego, CA, USA, June 13 - 17, 2022, pages 152–167. ACM, 2022.  doi:10.1145/3519939.3523711.  
 
\par\noindent  [22]  Serdar Erbatur, Andrew M. Marshall, and Christophe Ringeissen.  Non-disjoint combined unification and closure by equational paramodulation.  In Boris Konev and Giles Reger, editors, Frontiers of Combining Systems - 13th International Symposium, FroCoS 2021, Birmingham, UK, September 8-10, 2021, Proceedings, volume 12941 of Lecture Notes in Computer Science, pages 25–42. Springer, 2021.  doi:10.1007/978-3-030-86205-3\_2.  
 
\par\noindent  [23]  Serdar Erbatur, Andrew M. Marshall, and Christophe Ringeissen.  Combined Hierarchical Matching: the Regular Case.  In Amy P. Felty, editor, 7th International Conference on Formal Structures for Computation and Deduction, FSCD 2022, August 2-5, 2022, Haifa, Israel, volume 228 of LIPIcs, pages 6:1–6:22. Schloss Dagstuhl – Leibniz-Zentrum für Informatik, 2022.  doi:10.4230/LIPICS.FSCD.2022.6.  
 
\par\noindent  [24]  Boris Galitsky.  Developing Enterprise Chatbots - Learning Linguistic Structures.  Springer, 2019.  doi:10.1007/978-3-030-04299-8.  
 
\par\noindent  [25]  Peter Graf.  Substitution tree indexing.  In Jieh Hsiang, editor, Rewriting Techniques and Applications, 6th International Conference, RTA-95, Kaiserslautern, Germany, April 5-7, 1995, Proceedings, volume 914 of Lecture Notes in Computer Science, pages 117–131. Springer, 1995.  doi:10.1007/3-540-59200-8\_52.  
 
\par\noindent  [26]  Alexander Herold.  Combination of unification algorithms in equational theories.  PhD thesis, Kaiserslautern University of Technology, Germany, 1987.  URL: https://d-nb.info/880327979.  
 
\par\noindent  [27]  Jean-Pierre Jouannaud.  Syntactic theories.  In Branislav Rovan, editor, Mathematical Foundations of Computer Science 1990, MFCS’90, Banská Bystrica, Czechoslovakia, August 27-31, 1990, Proceedings, volume 452 of Lecture Notes in Computer Science, pages 15–25. Springer, 1990.  doi:10.1007/BFB0029593.  
 
\par\noindent  [28]  Claude Kirchner.  Méthodes et outils de conception systématique d’algorithmes d’unification dans les théories équationnelles.  Thèses d’état, Université de Nancy I, 1985.  
 
\par\noindent  [29]  Claude Kirchner and Francis Klay.  Syntactic theories and unification.  In Proceedings of the Fifth Annual Symposium on Logic in Computer Science (LICS ’90), Philadelphia, Pennsylvania, USA, June 4-7, 1990, pages 270–277. IEEE Computer Society, 1990.  doi:10.1109/LICS.1990.113753.  
 
\par\noindent  [30]  Hélène Kirchner and Christophe Ringeissen.  Combining symbolic constraint solvers on algebraic domains.  J. Symb. Comput., 18(2):113–155, 1994.  doi:10.1006/JSCO.1994.1040.  
 
\par\noindent  [31]  Boris Konev and Temur Kutsia.  Anti-unification of concepts in description logic EL.  In Chitta Baral, James P. Delgrande, and Frank Wolter, editors, Principles of Knowledge Representation and Reasoning: Proceedings of the Fifteenth International Conference, KR 2016, Cape Town, South Africa, April 25-29, 2016, pages 227–236. AAAI Press, 2016.  URL: http://www.aaai.org/ocs/index.php/KR/KR16/paper/view/12880.  
 
\par\noindent  [32]  Christopher Lynch and Barbara Morawska.  Basic syntactic mutation.  In Andrei Voronkov, editor, Automated Deduction - CADE-18, 18th International Conference on Automated Deduction, Copenhagen, Denmark, July 27-30, 2002, Proceedings, volume 2392 of Lecture Notes in Computer Science, pages 471–485. Springer, 2002.  doi:10.1007/3-540-45620-1\_37.  
 
\par\noindent  [33]  Sonu Mehta, Ranjita Bhagwan, Rahul Kumar, Chetan Bansal, Chandra Shekhar Maddila, B. Ashok, Sumit Asthana, Christian Bird, and Aditya Kumar.  Rex: Preventing bugs and misconfiguration in large services using correlated change analysis.  In Ranjita Bhagwan and George Porter, editors, 17th USENIX Symposium on Networked Systems Design and Implementation, NSDI 2020, Santa Clara, CA, USA, February 25-27, 2020, pages 435–448. USENIX Association, 2020.  URL: https://www.usenix.org/conference/nsdi20/presentation/mehta.  
 
\par\noindent  [34]  Tobias Nipkow.  Proof transformations for equational theories.  In Proceedings of the Fifth Annual Symposium on Logic in Computer Science (LICS ’90), Philadelphia, Pennsylvania, USA, June 4-7, 1990, pages 278–288. IEEE Computer Society, 1990.  doi:10.1109/LICS.1990.113754.  
 
\par\noindent  [35]  Tobias Nipkow.  Combining matching algorithms: The regular case.  J. Symb. Comput., 12(6):633–654, 1991.  doi:10.1016/S0747-7171(08)80145-9.  
 
\par\noindent  [36]  Julian Parsert and Elizabeth Polgreen.  Reinforcement learning and data-generation for syntax-guided synthesis.  In Michael J. Wooldridge, Jennifer G. Dy, and Sriraam Natarajan, editors, Thirty-Eighth AAAI Conference on Artificial Intelligence, AAAI 2024, Thirty-Sixth Conference on Innovative Applications of Artificial Intelligence, IAAI 2024, Fourteenth Symposium on Educational Advances in Artificial Intelligence, EAAI 2014, February 20-27, 2024, Vancouver, Canada, pages 10670–10678. AAAI Press, 2024.  doi:10.1609/AAAI.V38I9.28938.  
 
\par\noindent  [37]  Brigitte Pientka.  Higher-order term indexing using substitution trees.  ACM Trans. Comput. Log., 11(1):6:1–6:40, 2009.  doi:10.1145/1614431.1614437.  
 
\par\noindent  [38]  Gordon D. Plotkin.  A note on inductive generalization.  Machine Intel., 5(1):153–163, 1970.  
 
\par\noindent  [39]  Stanislaw J. Purgal, David M. Cerna, and Cezary Kaliszyk.  Learning higher-order logic programs from failures.  In Luc De Raedt, editor, Proceedings of the Thirty-First International Joint Conference on Artificial Intelligence, IJCAI 2022, Vienna, Austria, 23-29 July 2022, pages 2726–2733. ijcai.org, 2022.  doi:10.24963/IJCAI.2022/378.  
 
\par\noindent  [40]  Mohammad Raza, Sumit Gulwani, and Natasa Milic-Frayling.  Programming by Example Using Least General Generalizations.  In Carla E. Brodley and Peter Stone, editors, Proceedings of the Twenty-Eighth AAAI Conference on Artificial Intelligence, July 27 -31, 2014, Québec City, Québec, Canada, pages 283–290. AAAI Press, 2014.  doi:10.1609/AAAI.V28I1.8744.  
 
\par\noindent  [41]  John C. Reynolds.  Transformational systems and the algebraic structure of atomic formulas.  Machine Intel., 5(1):135–151, 1970.  
 
\par\noindent  [42]  Manfred Schmidt-Schauß.  Unification in a combination of arbitrary disjoint equational theories.  J. Symb. Comput., 8(1/2):51–99, 1989.  doi:10.1016/S0747-7171(89)80022-7.  
 
\par\noindent  [43]  Erik Tidén.  Unification in combinations of collapse-free theories with disjoint sets of function symbols.  In Jörg H. Siekmann, editor, 8th International Conference on Automated Deduction, Oxford, England, July 27 - August 1, 1986, Proceedings, volume 230 of Lecture Notes in Computer Science, pages 431–449. Springer, 1986.  doi:10.1007/3-540-16780-3\_110.  
 
\par\noindent  [44]  Emily Rowan Winter, Vesna Nowack, David Bowes, Steve Counsell, Tracy Hall, Sæmundur Óskar Haraldsson, John R. Woodward, Serkan Kirbas, Etienne Windels, Olayori McBello, Abdurahman Atakishiyev, Kevin Kells, and Matthew W. Pagano.  Towards developer-centered automatic program repair: findings from Bloomberg.  In Abhik Roychoudhury, Cristian Cadar, and Miryung Kim, editors, Proceedings of the 30th ACM Joint European Software Engineering Conference and Symposium on the Foundations of Software Engineering, ESEC/FSE, Singapore, November 14-18, 2022, pages 1578–1588. ACM, 2022.  doi:10.1145/3540250.3558953.  
 
\par\noindent  [45]  Kaiyu Yang and Jia Deng.  Learning symbolic rules for reasoning in quasi-natural language.  Trans. Mach. Learn. Res., 2023, 2023.  URL: https://openreview.net/forum?id=gwRwHUZUgz.  
 
\par\noindent  [46]  Katherine A. Yelick.  Unification in combinations of collapse-free regular theories.  J. Symb. Comput., 3(1/2):153–181, 1987.  doi:10.1016/S0747-7171(87)80025-1.  
 
\par\noindent  [47]  Gonzague Yernaux and Wim Vanhoof.  Anti-unification of unordered goals.  In Florin Manea and Alex Simpson, editors, 30th EACSL Annual Conference on Computer Science Logic, CSL 2022, February 14-19, 2022, Göttingen, Germany (Virtual Conference), volume 216 of LIPIcs, pages 37:1–37:17. Schloss Dagstuhl – Leibniz-Zentrum für Informatik, 2022.  doi:10.4230/LIPICS.CSL.2022.37.  
 
\par\noindent  [48]  Gonzague Yernaux and Wim Vanhoof.  On detecting semantic clones in constraint logic programs.  In 16th IEEE International Workshop on Software Clones, IWSC 2022, Limassol, Cyprus, October 2, 2022, pages 32–38. IEEE, 2022.  doi:10.1109/IWSC55060.2022.00014.  
 
\par\noindent  [49]  Dongjun Youn, Wonho Shin, Jaehyun Lee, Sukyoung Ryu, Joachim Breitner, Philippa Gardner, Sam Lindley, Matija Pretnar, Xiaojia Rao, Conrad Watt, and Andreas Rossberg.  Bringing the WebAssembly Standard up to Speed with SpecTec.  Proc. ACM Program. Lang., 8(PLDI):1559–1584, 2024.  doi:10.1145/3656440.  
  
% END bib
